% !PRINTED 211
\weberpage{211}
Durch Transformation mittelst der Substitution

\[
\begin{pmatrix} e^{\frac{\pi i}{4}}, & 0 \\ 0, & e^{-\frac{\pi i}{4}} \end{pmatrix}
\]

geht die Di\"edergruppe $D_m$ in

\[
\begin{pmatrix} \varepsilon^{1/2\,\lambda}, & 0 \\ 0, & \varepsilon^{-1/2\,\lambda} \end{pmatrix}, \quad
\begin{pmatrix} 0, & \varepsilon^{1/2\,\lambda} \\ -\varepsilon^{-1/2\,\lambda}, & 0 \end{pmatrix}
\]

oder in

\[
\begin{gathered}
x, \; \varepsilon x, \; \varepsilon^2 x, \; \ldots, \; \varepsilon^{m-1} x\\[4pt]
\dfrac{-1}{x}, \; \dfrac{-\varepsilon}{x}, \; \dfrac{-\varepsilon^2}{x}, \; \cdots, \; \dfrac{-\varepsilon^{m-1}}{x}
\end{gathered}
\]

\"uber, und die Grundformen f\"ur diese Darstellung sind:

\[
f_1 = x_1 x_2, \quad f_2 = x_1^m + i^m x_2^m, \quad f_3 = x_1^m - i^m x_2^m.
\]

Die Di\"edergruppe $D_m$ enth\"alt als Theiler die cyklische
Gruppe $C_m$, und zwar ist $C_m$ ein \gesperrt{Normaltheiler} von $D_m$.

Durch Transformation mit irgend einer multiplicativen
Substitution

\[
\begin{pmatrix} \alpha, & 0 \\ 0, & \delta \end{pmatrix}
\]

\"andert sich $C_m$ nicht, w\"ahrend die nicht in $C_m$ enthaltenen Substitutionen
 von $D_m$ ihre Form \"andern; denn es ist:

\[
\begin{gathered}
\begin{pmatrix} \delta, & 0 \\ 0, & \alpha \end{pmatrix}
\begin{pmatrix} \varepsilon, & 0 \\ 0, & 1 \end{pmatrix}
\begin{pmatrix} \alpha, & 0 \\ 0, & \delta \end{pmatrix}
=
\begin{pmatrix} \varepsilon, & 0 \\ 0, & 1 \end{pmatrix},\\[6pt]
\begin{pmatrix} \delta, & 0 \\ 0, & \alpha \end{pmatrix}
\begin{pmatrix} 0, & \varepsilon \\ 1, & 0 \end{pmatrix}
\begin{pmatrix} \alpha, & 0 \\ 0, & \delta \end{pmatrix}
=
\begin{pmatrix} 0, & \varepsilon \delta^2 \\ \alpha^2, & 0 \end{pmatrix}.
\end{gathered}
\]

Abgesehen von einer solchen multiplicativen Substitution ist
aber die Gruppe $D_m$ durch die darin enthaltene Gruppe $C_m$
vollkommen bestimmt. Denn setzen wir

\[
D_m = C_m + \varphi C_m,
\]
worin

\[
\varphi = \begin{pmatrix} \alpha, & \beta \\ \gamma, & \delta \end{pmatrix}
\]

gesetzt ist, so muss, da $C_m$ Normaltheiler von $D_m$ sein muss,
$\varphi C_m = C_m \varphi$ sein, d.\,h. es muss sich zu jedem Exponenten $r$
ein Exponent $s$, und umgekehrt, bestimmen lassen, so dass

\[
\begin{pmatrix} \alpha, & \beta \\ \gamma, & \delta \end{pmatrix}
\begin{pmatrix} \varepsilon^r, & 0 \\ 0, & 1 \end{pmatrix}
=
\begin{pmatrix} \varepsilon^s, & 0 \\ 0, & 1 \end{pmatrix}
\begin{pmatrix} \alpha, & \beta \\ \gamma, & \delta \end{pmatrix}
\]
%CONT>>
